\subsection*{本章练习与挑战解析}
\addcontentsline{toc}{subsection}{本章练习与挑战解析}

\Needspace{8\baselineskip}

\block{条件求值：把已知整体用起来}

\Needspace{7\baselineskip}

\begin{solution}[before skip=4pt plus 1pt,after skip=4pt plus 1pt,top=6pt,bottom=6pt]{练习 16 · 解析}

\think 左边加上 $2a^2b^2$ 后成为 $(a^2+b^2)^2$，右边也只含 $a^2+b^2$。这时可把目标整体设为一个新字母。

\solve 移项得 $(a^2+b^2)^2=a^2+b^2+6$。令 $s=a^2+b^2\ge0$，则
\[s^2-s-6=(s-3)(s+2)=0.\]
所以 $s=3$ 或 $s=-2$。由 $s\ge0$ 排除 $-2$，故 $a^2+b^2=3$。

\insight{换元后保留整体自带的范围条件，才能排除不符合原题的解。}

\end{solution}

\Needspace{7\baselineskip}

\begin{solution}[before skip=4pt plus 1pt,after skip=4pt plus 1pt,top=6pt,bottom=6pt]{练习 17 · 解析}

\think 取值相等可以转为差为零。分解后，一个因式是 $m-n$，题目已经说明它不为零，因此另一个因式必须为零。

\solve 两式作差得 $(m-n)(m+n-2)=0$。由 $m\ne n$ 得 $m+n=2$。
于是
\[4=(m+n)^2=m^2+n^2+2mn=5+2mn,\]
所以 $mn=-\frac12$。又
\[a=m(m-2)=m(-n)=-mn=\frac12.\]

\insight{取值相等先作差；非零条件往往用来排除其中一个零因式。}

\end{solution}

\Needspace{8\baselineskip}

\block{简便计算：先分解，再约分}

\Needspace{7\baselineskip}

\begin{solution}[before skip=4pt plus 1pt,after skip=4pt plus 1pt,top=6pt,bottom=6pt]{练习 18 · 解析}

\think 方法拓展章中的添项思想仍可使用。$t^4+64=(t^2+8)^2-(4t)^2$。用 $F(t)$ 表示代数式 $t^2+4t+8$，另一个因式正好是 $F(t-4)$，数字间相差 $4$ 就成为约分依据。

\solve 设 $F(t)=t^2+4t+8$，则
\[t^4+64=(t^2+4t+8)(t^2-4t+8)=F(t)F(t-4).\]
分子各项分解后为
\[F(7)F(3)\,F(15)F(11)\,F(23)F(19)\,F(31)F(27)\,F(39)F(35),\]
分母各项分解后为
\[F(3)F(-1)\,F(11)F(7)\,F(19)F(15)\,F(27)F(23)\,F(35)F(31).\]
因为 $F(t)=(t+2)^2+4>0$，可以约去相同因式，得到
\[\text{原式}=\frac{F(39)}{F(-1)}=\frac{1685}{5}=337.\]

\insight{把两个二次因式写成同一个式子在不同数处的取值，可看清错位约分。}

\end{solution}

\Needspace{8\baselineskip}

\block{平方与非负性：把条件变成限制}

\Needspace{7\baselineskip}

\begin{solution}[before skip=4pt plus 1pt,after skip=4pt plus 1pt,top=6pt,bottom=6pt]{练习 19 · 解析}

\think 为了产生 $12xy$，先用 $(x+6y)^2$，它用掉 $36y^2$；剩下 $16y^2-8y+1$ 又是一个完全平方。

\solve \[x^2+12xy+52y^2-8y+1=(x+6y)^2+(4y-1)^2=0.\]
两个平方均非负，因此 $x+6y=0$、$4y-1=0$，得 $y=\frac14$、$x=-\frac32$。
所以 $x^2-y^2=\frac94-\frac1{16}=\frac{35}{16}$。

\insight{先根据交叉项确定第一个平方，再检查余下的项是否能配成第二个平方。}

\end{solution}

\Needspace{7\baselineskip}

\begin{solution}[before skip=4pt plus 1pt,after skip=4pt plus 1pt,top=6pt,bottom=6pt]{练习 20 · 解析}

\think 两次作差得到同一个二次式 $m^2-mn+n^2$。把它写成平方和，就能判断严格正负，不需要判别式。

\solve \[a-c=m^2-mn+n^2,\qquad c-b=m^2-mn+n^2.\]
而
\[m^2-mn+n^2=\frac12[(m-n)^2+m^2+n^2]>0,\]
其中严格大于零由 $m\ne n$ 保证。因此 $a>c$、$c>b$，即 $b<c<a$。

\insight{比较大小先作差；严格不等式还要说明平方和为何不可能为零。}

\end{solution}

\Needspace{7\baselineskip}

\begin{solution}[before skip=4pt plus 1pt,after skip=4pt plus 1pt,top=6pt,bottom=6pt]{练习 21 · 解析}

\think 配方后两个平方的和不超过 $2$，每个整数的绝对值至多为 $1$。两坐标的取值范围因此很小，可以直接完整列举。

\solve 原不等式等价于 $(x-1)^2+(y-1)^2\le2$。
因为 $x,y$ 是整数，$x-1$、$y-1$ 各只能为 $-1,0,1$，即 $x,y\in\{0,1,2\}$。
这 $3\times3=9$ 种组合的平方和都不超过 $2$，故共有 $9$ 个整数点。

\insight{先用非负性缩小取值范围，再列举全部组合；不必借助圆的知识。}

\end{solution}

\Needspace{8\baselineskip}

\block{整数问题：把等式转成有限因数对}

\Needspace{7\baselineskip}

\begin{solution}[before skip=4pt plus 1pt,after skip=4pt plus 1pt,top=6pt,bottom=6pt]{练习 22 · 解析}

\think 展开该乘积得到 $6ab-9a+10b-15$，利用已知条件可得到积为 $288$。新因式必须分别满足除以 $3$ 的余数和奇偶限制。

\solve 原式等价于 $(3a+5)(2b-3)=288$。记 $p=3a+5$、$q=2b-3$，则 $p-5$ 必须是 $3$ 的倍数，$q$ 为奇数。
因为 $288=2^5\times3^2$，$q$ 只能取 $\pm1,\pm3,\pm9$。
对应的 $p=288/q$ 为 $288,96,32,-288,-96,-32$。
其中只有 $p=32$ 使 $p-5$ 是 $3$ 的倍数，故 $p=32$、$q=9$。
于是 $a=9$、$b=6$，$a+b=15$。

\insight{整数因数对包含负因数；系数带来的余数限制可以有效筛选。}

\end{solution}

\Needspace{7\baselineskip}

\begin{solution}[before skip=4pt plus 1pt,after skip=4pt plus 1pt,top=6pt,bottom=6pt]{练习 23 · 解析}

\think 总和为 $703$。剩余数之和是 $703-a-b$，等于 $ab$ 后可通过加 $1$ 写成 $(a+1)(b+1)$，再用数的范围筛选。

\solve 设所选两数为 $a<b$，则 $1\le a<b\le37$，且
\[703-a-b=ab,\qquad (a+1)(b+1)=704.\]
令 $u=a+1$、$v=b+1$，则 $2\le u<v\le38$。
因 $uv=704$，有 $u\ge\frac{704}{38}>18$，且 $u<\sqrt{704}<27$。
$704=2^6\times11$，在整数 $19$ 至 $26$ 中，只有 $22$ 是其因数。
故 $u=22$、$v=32$，即 $a=21$、$b=31$，两数之差为 $10$。

\insight{列因数对之前先使用范围限制，能减少枚举并保证不漏解。}

\end{solution}

\Needspace{7\baselineskip}

\begin{solution}[before skip=4pt plus 1pt,after skip=4pt plus 1pt,top=6pt,bottom=6pt]{练习 24 · 解析}

\think 移项配方可得到 $(a+1)^2-b^2=2019$。因为 $a+1+b>0$，乘积为正迫使另一因数也为正，只需列正因数对。

\solve 原等式等价于 $(a+1-b)(a+1+b)=2019=3\times673$。
两个因数均为正整数，且前者小于后者，只有 $(1,2019)$ 和 $(3,673)$ 两对。
第一对给出 $a=b=1009$，不符合 $a\ne b$；第二对给出 $a=337$、$b=335$。
所以 $a+b=672$。

\insight{配方把平方差转为整数乘积，再逐一检查正数与不相等条件。}

\end{solution}

\Needspace{8\baselineskip}

\block{数的性质：整除、质数与完全平方}

\Needspace{7\baselineskip}

\begin{solution}[before skip=4pt plus 1pt,after skip=4pt plus 1pt,top=6pt,bottom=6pt]{练习 25 · 解析}

\think 添减 $36n^2$ 后可以用平方差。得到两个正整数因数后，质数条件要求其中一个为 $1$，还要检查最后所得数确为质数。

\solve \begin{align*}
n^4-16n^2+100&=(n^2+10)^2-(6n)^2\\
&=(n^2+6n+10)(n^2-6n+10).
\end{align*}
第一因数 $n^2+6n+10\ge17>1$；第二因数为 $(n-3)^2+1\ge1$。
因此质数条件迫使 $(n-3)^2+1=1$，得 $n=3$。
代回值为 $37$。$37$ 不被 $2,3,5$ 整除，且 $\sqrt{37}<7$，所以它是质数。

\insight{“质数使一个因数为一”必须先说明两个因数都是正整数。}

\end{solution}

\Needspace{7\baselineskip}

\begin{solution}[before skip=4pt plus 1pt,after skip=4pt plus 1pt,top=6pt,bottom=6pt]{练习 26 · 解析}

\think 两数相邻。按 $t$ 展开后，系数依次为 $1,2,3,2,1$，与方法拓展章的完全平方结构一致。

\solve 令 $t=2018$，则
\begin{align*}
t^2+t^2(t+1)^2+(t+1)^2&=t^4+2t^3+3t^2+2t+1\\
&=(t^2+t+1)^2=4074343^2.
\end{align*}
所以原数是完全平方数。又 $t(t+1)$ 为偶数，$t^2+t+1$ 为奇数，所以原数也是奇数。

\insight{相邻数先写成 $t,t+1$，数的规律就变成多项式的结构。}

\end{solution}

\Needspace{8\baselineskip}

\block{几何中的数量关系}

\Needspace{6\baselineskip}

\begin{solution}[before skip=4pt plus 1pt,after skip=4pt plus 1pt,top=6pt,bottom=6pt]{练习 27 · 解析}

\think 等式中同时含三个平方与三个两两乘积。将等式两边乘以 $2$，可以把各交叉项配入 $(a-b)^2$、$(b-c)^2$、$(c-a)^2$。

\solve 由已知条件得

\begin{align*}
2a^2+2b^2+2c^2-2ab-2bc-2ca&=0,\\
(a-b)^2+(b-c)^2+(c-a)^2&=0.
\end{align*}

三个平方均非负，和为零，故 $a-b=b-c=c-a=0$，即 $a=b=c$。因此 $\triangle ABC$ 为等边三角形。

同一变形还说明，对任意实数 $a,b,c$，有

\[
a^2+b^2+c^2-ab-bc-ca
=\frac12[(a-b)^2+(b-c)^2+(c-a)^2]\ge0,
\]

等号成立当且仅当 $a=b=c$。

\insight{平方和为零可以逐项推出零，是因为每一项都非负；再把字母相等翻译成图形结论。}

\end{solution}

\Needspace{8\baselineskip}

\block{选学挑战}

\Needspace{7\baselineskip}

\begin{solution}[before skip=4pt plus 1pt,after skip=4pt plus 1pt,top=6pt,bottom=6pt]{挑战 2 · 解析}

\think 新数加 $1$ 恰好等于两个旧数各加 $1$ 后的乘积。因而所有在场的数各加 $1$ 后再相乘，这个总乘积在每次操作中都保持不变。

\solve 设一次操作写入 $c=a+b+ab$，则 $c+1=(a+1)(b+1)$。
初始时所有数各加 $1$ 后的乘积为
\[2\times\frac32\times\frac43\times\cdots\times\frac{101}{100}=101.\]
每次操作保持这个乘积不变。设最终剩下 $r$，则 $r+1=101$，故 $r=100$。

\insight{操作式含 $a+b+ab$ 时，先看“各加一之后的乘积”是否保持不变。}

\end{solution}

\Needspace{7\baselineskip}

\begin{solution}[before skip=4pt plus 1pt,after skip=4pt plus 1pt,top=6pt,bottom=6pt]{挑战 3 · 解析}

\think 作差后可以消去非零因式 $a-b$，将 $2019$ 表示成关于 $a,b$ 的积。代入 $c$ 后，四次项抵消，留下能直接判断符号的因式。

\solve 由 $a^4-b^4=2019(a-b)$，且 $a\ne b$，得
\[(a+b)(a^2+b^2)=2019.\]
所以 $a+b>0$，并且
\begin{align*}
c&=a^4-a(a+b)(a^2+b^2)\\
&=-ab(a^2+ab+b^2).
\end{align*}
因为 $a^2+ab+b^2=\frac12[(a+b)^2+a^2+b^2]>0$，而 $c>0$，故 $ab<0$。
再由
\[c-(ab)^2=-ab(a+b)^2>0,\]
得 $(ab)^2<c$。结合 $ab<0$，得到 $-\sqrt c<ab<0$。

\insight{同值条件先作差；把得到的低次关系代回原式，能把高次问题转成符号判断。}

\end{solution}
